\subsubsection{Boundary Integral Methods}\label{subsec:boundary-integral-methods}

Let $\Omega$ be a bounded domain with a smooth boundary. We want to solve the problem
\begin{equation}\label{eq:harmonic-dirichlet-boundary-integral-problem}
\Delta u=0,\ \bm{x}\in\Omega,\qquad u=\phi(\bm{x}),\ \bm{x}\in\partial\Omega.
\end{equation}
If we knew the Green\textquotesingle s function, we would have the representation
\begin{equation}\label{eq:harmonic-green-normal-boundary-representation}
u(\bm{x})=\int_{\partial\Omega}\dfrac{\partial}{\partial n_{\bm{y}}}G(\bm{x},\bm{y})\phi(\bm{y})\,\mathrm{d}S_{\bm{y}}.
\end{equation}
We make an ansatz analogous to \eqref{eq:harmonic-green-normal-boundary-representation}, with the Green\textquotesingle s function replaced by the fundamental solution of the Laplace equation
\begin{equation}\label{eq:laplace-double-layer-potential-ansatz}
u(\bm{x})=\int_{\partial\Omega}\dfrac{\partial}{\partial n_{\bm{y}}}\mathcal{G}(|\bm{x}-\bm{y}|)g(\bm{y})\,\mathrm{d}S_{\bm{y}}.
\end{equation}
Here the function $g$ is unknown, and we are seeking an equation relating $g$ to $\phi$.

We note that for any $g\in C(\partial\Omega)$, the function $u$ given by \eqref{eq:laplace-double-layer-potential-ansatz} is harmonic in $\Omega$; we can simply take the Laplacian with respect to $\bm{x}$ under the integral. To satisfy the boundary condition, we must have
\begin{equation}\label{eq:double-layer-dirichlet-boundary-limit}
\phi(\bm{x})=\lim_{\bm{z}\to\bm{x},\bm{z}\in\Omega}\int_{\partial\Omega}\dfrac{\partial}{\partial n_{\bm{y}}}\mathcal{G}(|\bm{z}-\bm{y}|)g(\bm{y})\,\mathrm{d}S_{\bm{y}}
\end{equation}
for $\bm{x}\in\partial\Omega$; this is the desired equation relating $g$ to $\phi$. One cannot pass to the limit in \eqref{eq:double-layer-dirichlet-boundary-limit} by simply substituting $\bm{x}$ for $\bm{z}$; although the integral exists for $\bm{z}\in\partial\Omega$, it is discontinuous there. Indeed, we shall show below that actually
\begin{equation}\label{eq:double-layer-interior-jump-relation}
\begin{aligned}
&\lim_{\bm{z}\to\bm{x},\bm{z}\in\Omega}\int_{\partial\Omega}\dfrac{\partial}{\partial n_{\bm{y}}}\mathcal{G}(|\bm{z}-\bm{y}|)g(\bm{y})\,\mathrm{d}S_{\bm{y}}\\
&\quad=\int_{\partial\Omega}\dfrac{\partial}{\partial n_{\bm{y}}}\mathcal{G}(|\bm{x}-\bm{y}|)g(\bm{y})\,\mathrm{d}S_{\bm{y}}+\dfrac12g(\bm{x}).
\end{aligned}
\end{equation}


Recall that a similar situation applies to Cauchy\textquotesingle s formula; if $C$ is a smooth closed curve in the plane and $f$ is analytic, then
\begin{equation}\label{eq:cauchy-integral-boundary-principal-value-analogue}
\dfrac1{2\pi i}\int_C\dfrac{f(\zeta)}{\zeta-z}\,\mathrm{d}\zeta
\end{equation}
equals $f(z)$ for $z$ inside $C$, 0 for $z$ outside $C$ and (in the sense of principal value) $f(z)/2$ on $C$. Inserting \eqref{eq:double-layer-interior-jump-relation} in \eqref{eq:double-layer-dirichlet-boundary-limit}, we obtain
\begin{equation}\label{eq:dirichlet-double-layer-boundary-integral-equation}
\phi(\bm{x})-\dfrac12g(\bm{x})=\int_{\partial\Omega}\dfrac{\partial}{\partial n_{\bm{y}}}\mathcal{G}(|\bm{x}-\bm{y}|)g(\bm{y})\,\mathrm{d}S_{\bm{y}}.
\end{equation}


We have thus replaced the partial differential equation \eqref{eq:harmonic-dirichlet-boundary-integral-problem} by the equivalent integral equation \eqref{eq:dirichlet-double-layer-boundary-integral-equation}. This has two advantages. As we shall see in Chapter~\ref{chap:linear-elliptic-equations}, it is fairly easy to develop an existence theory for integral equations such as \eqref{eq:dirichlet-double-layer-boundary-integral-equation}. Moreover, a numerical approach based on \eqref{eq:dirichlet-double-layer-boundary-integral-equation} rather than \eqref{eq:harmonic-dirichlet-boundary-integral-problem} has the advantage of working with a problem in a lower space dimension, which translates into fewer gridpoints. Indeed, there is an extensive literature on ``boundary-element methods'' for Laplace\textquotesingle s equation as well as for the Stokes equation.

It remains to verify \eqref{eq:double-layer-interior-jump-relation}. Let $\bm{z}$ be close to $\partial\Omega$, and let $\bm{x}$ be the point on $\partial\Omega$ nearest to $\bm{z}$; without loss of generality, we may choose the coordinate system in such a way that $\bm{x}$ is the origin and the normal to $\partial\Omega$ is in the $m$th coordinate direction. Let $N$ be a neighborhood of the origin; we can then split up the right-hand side of \eqref{eq:double-layer-dirichlet-boundary-limit} as follows:
\begin{equation}\label{eq:double-layer-local-distant-splitting}
\begin{aligned}
\int_{\partial\Omega}\dfrac{\partial}{\partial n_{\bm{y}}}\mathcal{G}(|\bm{z}-\bm{y}|)g(\bm{y})\,\mathrm{d}S_{\bm{y}}
&=\int_{\partial\Omega\cap N}\dfrac{\partial}{\partial n_{\bm{y}}}\mathcal{G}(|\bm{z}-\bm{y}|)g(\bm{y})\,\mathrm{d}S_{\bm{y}}\\
&\quad+\int_{\partial\Omega\setminus N}\dfrac{\partial}{\partial n_{\bm{y}}}\mathcal{G}(|\bm{z}-\bm{y}|)g(\bm{y})\,\mathrm{d}S_{\bm{y}}.
\end{aligned}
\end{equation}


The second term is continuous at $\bm{z}=\bm{0}$. For the first term, we choose $N$ small enough so that $\partial\Omega\cap N$ can be represented in the form $y_m=\phi(y_1,\ldots,y_{m-1})$; we set $\bm{u}=(y_1,\ldots,y_{m-1})$. This leads to
\begin{equation}\label{eq:boundary-graph-normal-surface-measure}
\bm{n}_{\bm{y}}=(-\nabla\phi,1)/\sqrt{1+|\nabla\phi|^2},\quad\mathrm{d}S_{\bm{y}}=\sqrt{1+|\nabla\phi|^2}\mathrm{d}\bm{u}.
\end{equation}
We may choose $N$ in such a way that $\partial\Omega\cap N=\{(\bm{u},\phi(\bm{u}))\mid|\bm{u}|<\epsilon\}$. The first term on the right-hand side of \eqref{eq:double-layer-local-distant-splitting} now assumes the form
\begin{equation}\label{eq:double-layer-local-graph-kernel}
\begin{aligned}
&\int_{\{|\bm{u}|<\epsilon\}}\nabla_{\bm{y}}\mathcal{G}(|\bm{z}-(\bm{u},\phi(\bm{u}))|)\cdot(-\nabla\phi(\bm{u}),1)g(\bm{u},\phi(\bm{u}))\,\mathrm{d}\bm{u}\\
&\quad=\int_{\{|\bm{u}|<\epsilon\}}\dfrac{-\bm{u}\cdot\nabla\phi(\bm{u})+\phi(\bm{u})-z_m}
{\Omega_m(\sqrt{|\bm{u}|^2+|\phi(\bm{u})-z_m|^2})^m}g(\bm{u},\phi(\bm{u}))\,\mathrm{d}\bm{u}.
\end{aligned}
\end{equation}
We note that if we set $z_m=0$ in \eqref{eq:double-layer-local-graph-kernel}, then $-\bm{u}\cdot\nabla\phi(\bm{u})+\phi(\bm{u})$ is of order $|\bm{u}|^2$ as $\bm{u}\to\bm{0}$, hence the integrand is of order $|\bm{u}|^{-(m-2)}$, i.e., it is integrable. Although the integral exists for $\bm{z}=\bm{0}$, we cannot take the limit $z_m\to0$ under the integral. We shall now consider this limit with the constraint that $z_m<0$. The term which needs to be investigated is
\begin{equation}\label{eq:double-layer-singular-normal-kernel}
-z_m\int_{\{|\bm{u}|<\epsilon\}}\dfrac1{\Omega_m(\sqrt{|\bm{u}|^2+|\phi(\bm{u})-z_m|^2})^m}g(\bm{u},\phi(\bm{u}))\,\mathrm{d}\bm{u}.
\end{equation}


For small $|\bm{u}|$ and $|z_m|$, one has
\begin{equation}\label{eq:double-layer-graph-denominator-expansion}
\dfrac1{\sqrt{|\bm{u}|^2+(\phi(\bm{u})-z_m)^2}^{\,m}}=\dfrac1{\sqrt{|\bm{u}|^2+z_m^2}^{\,m}}(1+O(|z_m|)+O(|\bm{u}|^2)),
\end{equation}
and it is easily checked that only the leading contribution leads to a discontinuity in \eqref{eq:double-layer-singular-normal-kernel} as $z_m\to0$. It thus remains to consider the integral
\begin{equation}\label{eq:double-layer-leading-normal-kernel}
-z_m\int_{\{|\bm{u}|<\epsilon\}}\dfrac1{\Omega_m(\sqrt{|\bm{u}|^2+z_m^2})^m}g(\bm{u},\phi(\bm{u}))\,\mathrm{d}\bm{u}.
\end{equation}


We define
\begin{equation}\label{eq:double-layer-tangential-spherical-mean}
I_r(g)=\dfrac1{r^{m-2}\Omega_{m-1}}\int_{\{|\bm{u}|=r\}}g(\bm{u},\phi(\bm{u}))\,\mathrm{d}S.
\end{equation}
We substitute $\bm{u}=-z_m\bm{v}$ in \eqref{eq:double-layer-leading-normal-kernel}. This leads to the expression
\begin{equation}\label{eq:double-layer-scaled-radial-integral}
\begin{aligned}
&\int_{\{|\bm{v}|\leq-\epsilon/z_m\}}\dfrac1{\Omega_m(\sqrt{|\bm{v}|^2+1})^m}g(-z_m\bm{v},\phi(-z_m\bm{v}))\,\mathrm{d}\bm{v}\\
&\quad=\dfrac{\Omega_{m-1}}{\Omega_m}\int_0^{-\epsilon/z_m}\dfrac{r^{m-2}}{(r^2+1)^{m/2}}I_{-z_mr}(g)\,\mathrm{d}r.
\end{aligned}
\end{equation}
In the limit $z_m\to0-$, we obtain
\begin{equation}\label{eq:double-layer-half-density-jump-limit}
\dfrac{\Omega_{m-1}}{\Omega_m}g(\bm{0})\int_0^\infty\dfrac{r^{m-2}}{(r^2+1)^{m/2}}\,\mathrm{d}r=\dfrac12g(\bm{0}).
\end{equation}


Here we have used that
\begin{equation}\label{eq:double-layer-radial-gamma-integral}
\int_0^\infty\dfrac{r^{m-2}}{(r^2+1)^{m/2}}\,\mathrm{d}r=\dfrac{\Gamma(\frac{m-1}2)\sqrt\pi}{2\Gamma(\frac m2)}
\end{equation}
(see \cite[p.~292]{GR}) and the expression for $\Omega_m$ obtained in Problem~\ref{ex:unit-ball-volume-gamma-function}.
